\documentclass[12pt]{article}
\usepackage[utf8]{inputenc}
\usepackage[T1]{fontenc}
\usepackage{hyperref}
\title{Sudest Language Resources for K-12 Curriculum Development}
\author{Compiled by Claude Sonnet 5 (claude-sonnet-5) for Andr\'as Kornai}
\date{}
\begin{document}
\maketitle
\section{Introduction}
Sudest (ISO 639-3 tgo, Glottocode sude1239) is an Oceanic language of Papua New Guinea. It is spoken by about 2,000 people on Sudest (Tagula, Vanatinai) Island in Milne Bay Province. Speakers call it Vanatina. It belongs to the Nimoa--Sudest group of the Papuan Tip subgroup. Little material was found.
\section{Primary Educational Resources}
No primers, textbooks or school materials were found.
\section{Linguistic and Reference Materials}
\begin{itemize}
\item Wikipedia article on Sudest, covering dialects and classification. \url{https://en.wikipedia.org/wiki/Sudest_language}
\item Omniglot page on Sudest. \url{https://omniglot.com/writing/sudest.htm}
\end{itemize}
\section{Religious and Community Resources}
\begin{itemize}
\item Joshua Project language page for tgo. The searches indicated a Bible translation project begun in 1985, with a New Testament in progress around 2015. \url{https://www.joshuaproject.net/index.php/languages/tgo}
\item Joshua Project profile of the Tagula (Sud Est) people. \url{https://joshuaproject.net/people_groups/15179/PP}
\end{itemize}
\section{Digital and Audio Resources}
\begin{itemize}
\item Wikitongues page for tgo. \url{https://wikitongues.org/languages/tgo/}
\item Verbix page for Sudest. \url{https://docs.verbix.com/Languages/Sudest}
\end{itemize}
\section{Curriculum Development Context}
No school programme or published teaching material was found. Any curriculum would have to be built from scratch, with the Bible translation as the main known written source.
\end{document}